
\chapter{Association Analysis}

\index{association analysis}
\index{machine learning (ML)}

Classification isn't the only machine learning task. Another one is called
\textbf{association analysis}, in which the goal isn't to predict the class of
new data items, but rather to determine which features of a data set have
``interesting'' relationships. Usually, ``interesting'' means ``correlated in a
way that can be taken advantage of.''

The principles of association analysis are used in fields as diverse as
bioinformatics, recommender systems, and cybersecurity intrusion detection.
Another common application is \textbf{market basket analysis}, in which it can
be employed to detect which sets of products tend to be bought together. This
information is useful, for instance, to retailers who seek optimal product
placement and effective marketing campaigns.

\section{Market basket analysis}
\index{market basket analysis}
\index{VIP card}
\index{MVP card}
\index{transaction}

All competitive grocery stores today keep track of their customers' purchases
as they check out. Even if they don't have a rewards program to lure customers
into getting a store-specific ``VIP card,'' they still know which sets of
products were bought together in a single trip to the store. Each such trip,
consisting of multiple products being purchased, is called a
\textbf{transaction}.

This type of association analysis relies on a huge table (\texttt{DataFrame})
of 1's and 0's. Each row is a transaction, and each column is a product the
store sells. A ``1'' in an entry indicates that \textit{at least one} of that
product was bought in that transaction. (Interestingly, quantities typically
aren't stored; just binary information about ``did this purchase contain at
least one of that item or not?'') Here's one such data set:

\begin{center}
\begin{tabular}{|c|c|c|c|c|c|}
\hline
diet\_soda & apples & beer & chips\_ahoy & diapers & cat\_food \\
\hline
0 & 1 & 0 & 0 & 0 & 0 \\
1 & 1 & 0 & 0 & 0 & 1 \\
1 & 1 & 1 & 0 & 1 & 0 \\
0 & 0 & 1 & 1 & 1 & 0 \\
1 & 1 & 0 & 1 & 1 & 0 \\
1 & 1 & 1 & 1 & 1 & 1 \\
0 & 0 & 0 & 0 & 0 & 1 \\
. & . & . & . & . & . \\
\hline
\end{tabular}
\end{center}

The first customer in this set apparently just ran into the store to grab a bag
of apples, while the second one bought some Diet Coke, an apple, and a couple
bags of cat food, \textit{etc.}

\bigskip
Now, some terms:
\vspace{-.15in}

\begin{itemize}
\itemsep.1em

\index{itemset}

\item \textbf{itemset}: An itemset is a group of one or more items. For
instance, ``\{apples,beer,cat\_food,chips\_ahoy\}'' is an itemset, as is
``\{diet\_soda\}''.

\index{support count}

\item \textbf{support count}: The support count for a given itemset is
\textit{the fraction of transactions in the entire data set that contain that
itemset.} In the sample data above, the support count for ``\{diet\_soda\}'' is
$\frac{4}{7} = .571$, since four out of the seven transactions contained diet
soda. The support count for ``\{apples,chips\_ahoy,diapers\}'' is just
$\frac{2}{7} = .286$.

\index{association rule}
\index{rule, association}
\index{antecedent}
\index{consequent}

\item \textbf{association rule}: An association rule is a statement of a
relationship between two itemsets. It looks like this: ``\{A,B\} $\Rightarrow$
\{C\},'' which means ``people who buy A and B tend to also buy C.'' The itemset
before the arrow is called the \textbf{antecedent} and the one after it is
called the \textbf{consequent}. Association
rules are general correlations, not hard-and-fast mandates. Their reliability
can be measured with the metrics below:

\vspace{-.15in}

\begin{itemize}

    \index{support}

    \item \textbf{support}: An association rule's ``support\footnote{Note that
    the term \textbf{support} properly applies to a rule, whereas the term
    \textbf{support \textit{count}} applies to an itemset. The \texttt{mlxtend}
    package calls them both ``support,'' though.}'' is the fraction
    of transactions in the data set that contain \textit{both} its antecedent
    \textit{and} its consequent. It's a measure of how widespread and
    applicable it is. The rule ``\{diet\_soda, apples\} $\Rightarrow$
    \{beer\}'' has a support of just $\frac{2}{7} = .286$ in the data set
    above. A rule's support is always between 0 and 1.

    \index{confidence}

    \item \textbf{confidence}: An association rule's ``confidence,'' on the
    other hand, is a measure of its reliability. It's \textit{the fraction of
    the times that it could apply where it's actually correct}. In terms of
    support, it's $\frac{\textrm{support(whole
    rule)}}{\textrm{support\_count(antecedent)}}$ The rule
    ``\{diet\_soda,apples\} $\Rightarrow$ \{beer\}'' has a support of
    $\frac{2}{4} = .5$ in the data set above, since out of the four times that
    diet soda and apples were purchased together, beer was also purchased in
    two of them. A rule's confidence is always between 0 and 1.

    \index{lift}

    \item \textbf{lift}: Finally, an association rule's ``lift'' is a numerical
    measure of how much predictive power it has. The key concept is that we
    want to know how likely the consequent is to be satisfied when the
    antecedent is satisfied \textit{relative to how likely the consequent is to
    be satisfied anyway.} Quantitatively, it's $\frac{\textrm{confidence(whole
    rule)}}{\textrm{support\_count(consequent)}}$. A rule's lift is always
    \textit{between 1 and $\infty$.} A lift of 1 means it's essentially
    useless: knowing the antecedent is true doesn't really make the consequent
    any more likely. \end{itemize}

\end{itemize}

\vspace{-.1in}

If we crunch all the numbers, we want to look for association rules with high
support, high confidence, and high lift. Which of those three end up being
the greatest priority depends somewhat on the application. If a rule doesn't
have high support, it doesn't come up much, so it may not be very useful. If
it doesn't have high confidence, then the trend it signifies isn't super
reliable, and we need to be aware of that. And if it doesn't have high lift,
then its antecedent doesn't tell us a whole lot that we didn't already know.

\vspace{-.1in}

\section{Implementation: the \texttt{mlxtend} library}

\index{conda install@``\texttt{conda install}'' command}
\index{mlxtend@\texttt{mlxtend} (package)}
\index{scikit-learn@\texttt{scikit-learn} (package)}
\index{Raschka, Sebastian}

Surprisingly, \texttt{scikit-learn} doesn't seem to have an association
analysis implementation out of the box. (At least, Stephen couldn't find one.)
There is one in the open source \texttt{mlxtend} library, though, by Sebastian
Raschka \& friends. You may need to run ``\texttt{conda install mlxtend}'' at
the command line for Spyder to access it.

\index{synthetic data set}
\index{association}
\index{prior (probability)}

\smallskip
Let's take it for a spin. First, we'll generate a synthetic data set with a
known association. We'll stock five items for sale in our pretend store, each
of which will have a different prior probability of being purchased (about 60\%
of all transactions will have bananas, 80\% will have pizza, \textit{etc.})
However, we'll ensure that in addition to 2,000 purely random ones, we'll have
a solid \textit{thousand} transactions with both bacon and eggs. This should
guarantee that ``\{bacon\} $\Rightarrow$ \{eggs\}'' comes out to be a
meaningful association.\footnote{``...and \{eggs\} $\Rightarrow$ \{bacon\} too,
right?'' you ask. Well, not necessarily, or at least not as strong: since 70\%
of the random transactions have eggs anyway, and only 10\% of them have bacon,
we're going to get a lot of bacon-less egg-full transactions, which will dial
down the confidence of that association. But let's wait and see what
\texttt{mlxtend} says.} Figure~\ref{fig:synthGroceries}
(p.~\pageref{fig:synthGroceries}) gives some code to synthesize such a data set.

\begin{figure}[ht]
\centering
\index{banana (fruit)}
\index{coffee}
\index{bacon}
\index{egg}
\index{pizza}
\begin{Verbatim}[fontsize=\footnotesize,frame=single,framesep=3mm]
eggs = []
bacon = []
bananas = []
pizza = []
coffee = []

for i in range(2000):
    eggs.append(np.random.choice([1,0],p=[.7,.3]))
    bacon.append(np.random.choice([1,0],p=[.1,.9]))
    bananas.append(np.random.choice([1,0],p=[.6,.4]))
    pizza.append(np.random.choice([1,0],p=[.8,.2]))
    coffee.append(np.random.choice([1,0],p=[.3,.7]))

for i in range(1000):
    eggs.append(1)
    bacon.append(1)
    bananas.append(np.random.choice([1,0],p=[.6,.4]))
    pizza.append(np.random.choice([1,0],p=[.8,.2]))
    coffee.append(np.random.choice([1,0],p=[.3,.7]))

transactions = pd.DataFrame({ 'eggs':eggs, 'bacon':bacon,
    'bananas':bananas, 'pizza':pizza, 'coffee':coffee})
\end{Verbatim}
\caption{Synthesizing a grocery shopping cart data set.}
\label{fig:synthGroceries}
\end{figure}

\bigskip

Okay, let's analyze. First, let's shuffle the \texttt{DataFrame}'s rows so
they're in a random order (not strictly necessary, but just on principle) and
then let \texttt{mlxtend} do its work:

\index{sample@\texttt{.sample()} (Pandas)}
\index{import@\texttt{import}}

\begin{Verbatim}[fontsize=\small,samepage=true,frame=single,framesep=3mm]
transactions = transactions.sample(frac=1)

from mlxtend.frequent_patterns import apriori
from mlxtend.frequent_patterns import association_rules

itemsets = apriori(transactions, min_support=0.1,
    use_colnames=True)
\end{Verbatim}

(Note that passing \texttt{frac=1} to the \texttt{.sample()} method says ``give
me \textit{all} the rows, please; but in a shuffled order.)

\index{apriori algorithm}
\index{min\_support@\texttt{min\_support}}

This code uses the so-called \textbf{apriori algorithm}. It computes the
support count \textit{only} for those itemsets that have a certain minimum
support level (in this case, \texttt{min\_support=0.1}). This is for simple
combinatorics reasons: if you think about it, the number of possible itemsets
is rather large if we have many columns. In fact, if we have $N$ columns, there
are $2^N$ different itemsets!\footnote{Sample numbers: if we have 100 different
items for sale in our store -- certainly not a very large number at all --
we'll have 1,267,650,600,228,229,401,496,703,205,376 different itemsets to
consider. This is far, far more than the number of atoms in the universe.} The
apriori algorithm is smart about not evaluating them all, because it uses an
easy trick: in order for an itemset to have greater than (say) .1 support
count, \textit{all of its subsets} must also have at least that much support.
This lets apriori whittle down greatly the number of itemsets it needs to
consider.

\bigskip
Let's print out the ten itemsets with the greatest support count:
\begin{Verbatim}[fontsize=\small,samepage=true,frame=single,framesep=3mm]
itemsets.sort_values('support',ascending=False).head(10)
\end{Verbatim}
\vspace{-.2in}

\begin{Verbatim}[fontsize=\small,samepage=true,frame=leftline,framesep=5mm,framerule=1mm]
               support                itemsets
              0.816000                  [eggs]
              0.801667                 [pizza]
              0.650667           [eggs, pizza]
              0.593333               [bananas]
              0.486000         [bananas, eggs]
              0.475333        [bananas, pizza]
              0.401667                 [bacon]
              0.387667  [bananas, eggs, pizza]
              0.380667           [bacon, eggs]
              0.316000          [bacon, pizza]
\end{Verbatim}

Not surprisingly, the one-item itemsets feature prominently in the list, with
eggs on the very top since not only did all of our second batch of 1000
transactions have eggs, but about 70\% of the others did as well.

\medskip

Now, let's generate (and compute stats on) all possible association rules
involving those item sets:

\begin{Verbatim}[fontsize=\footnotesize,samepage=true,frame=single,framesep=3mm]
rules = association_rules(itemsets)
print(rules.sort_values('confidence',ascending=False).head(10))
\end{Verbatim}
\vspace{-.2in}

\begin{Verbatim}[fontsize=\footnotesize,samepage=true,frame=leftline,framesep=2mm,framerule=1mm]
            antecedents consequents   support confidence     lift
        (bacon, coffee)      (eggs)   0.11900    0.95454  1.18626
         (bacon, pizza)      (eggs)   0.29966    0.94631  1.17603
                (bacon)      (eggs)   0.38100    0.94619  1.17588
       (bacon, bananas)      (eggs)   0.22500    0.94537  1.17486
(bacon, bananas, pizza)      (eggs)   0.17566    0.94275  1.17160
        (bacon, coffee)     (pizza)   0.10133    0.81283  1.02760
               (coffee)     (pizza)   0.24933    0.80777  1.02120
         (eggs, coffee)     (pizza)   0.20000    0.80645  1.01953
                (pizza)      (eggs)   0.63633    0.80446  0.99975
               (coffee)      (eggs)   0.24800    0.80345  0.99849
\end{Verbatim}

As expected, ``\{bacon\} $\Rightarrow$ \{eggs\}'' has a very high confidence
(and support, and lift), as do rules of the form ``\{bacon + various other
stuff\} $\Rightarrow$ \{eggs\}.'' Even in this synthetic case, however, the
lift values aren't all that high numerically, so don't expect them to be.

\section{Applying the knowledge}

\index{bacon}
\index{egg}

This is the point where the ML ends and the domain application begins. What
should a store do if armed with these statistics? In this case, we've learned
that \{bacon\} $\Rightarrow$ \{eggs\} is a good bet. So we might ask ourselves:
what product(s) that we're currently not selling very frequently (at least to
bacon \& egg lovers) would we like to sell more of? That choice of product will
have lots to do with inventory levels, the amount of markup we have on the
product (\textit{i.e.}, how much profit we make on it), and so forth.

\index{bagel}

Let's say we've identified a target product to market to these bacon \& egg
lovers: a new brand of bagels that's targeted specifically to big breakfast
eaters. It might make sense to position this item in our store somewhere
\textit{between} the aisles that bacon and eggs are sold on, so that every
person buying both bacon and eggs will have to walk past it. This is just one
simple example, but profit margins in grocery stores are so small that the
difference between booming and going out of business can rely precisely on many
micro-decisions like this, each inspired by a micro-insight from the data.
